\section{Micro-level Assumptions}

In the following, I propose a new theoretical model that includes supply to explain the joint patterns of the spatial distribution of population and built environment and their dynamics. I will not consider exogenous amenities like waterfronts, elevation, or the central business district in this model. Such features could be integrated into the model in the next step.

Within this housing market model, two types of agents interact. The first type of agent is households that demand housing units. The second type is landlords, representing the market's supply side. I will introduce the assumed behavior of each type of agent in this section, motivating my assumptions with empirical literature. The following sections formalize these assumptions into a computational model and discuss the dynamics and results these assumptions create. 


\subsection{Demand: Preferences and Budget Constraints}

A specificity of the housing market is that higher incomes usually do not increase the number of housing units demanded by household \citep[353]{glaeserUrbanDeclineDurable2005}\footnote{At least for most households, higher incomes do not lead to more housing units being demanded by a household for own residency. Wealthy households might use housing as an asset, and some may have multiple residencies, but this is the exception. Housing, therefore, represents a \textit{unit demand} and housing units are \textit{substitutable goods}.}. I assume that every household only demands one housing unit, but they maximize the utility they get from its characteristics.

Inequality is a necessary condition for residential segregation. No spatial patterns could be observed without differences in the characteristics of a city's residents. However, income inequality also implies that households differ in their financial means to realize their housing preferences. Wealthy households can outcompete poorer households in housing units they find desirable and consequently have a large number of available options. Poor households, on the other hand, have minimal housing options due to affordability. Households can only choose among housing units that they can afford.

Of course, households are not always on the hunt for housing that marginally improves their utility, especially because moving itself is costly. More commonly, changes in the situation of households, like a job change, cohabitation, or childbirth, trigger a housing search. Especially for poor households, moves are responses to evictions and failing infrastructure, so households need to secure housing on short notice \citep{delucaNotJustLateral2020}. However, if households consider moving, what are their preferences? 

Given the two assumptions on households' preferences reviewed below, households choose between between the available housing units. Available units are unoccupied units or they can stay at their current location. If they move, it is either because they can improve their residential satisfaction or because rent of their current residence has surpassed their income disposable on housing. 

\subsubsection{Housing Quality}

There is a rich literature about which attributes of a housing unit households are willing to pay for, which implies a preference for them. Among the most important preferences of households are attributes of the housing unit themselves. These may include the floor size, number of rooms, quality of the structure, type of tenure, as well as facilities like pools, balconies, or lifts \citep{chauCriticalReviewLiterature2003, bayerEquilibriumModelSorting2004, bruecknerLecturesUrbanEconomics2011}. Although such preferences can be manifold and depend on the respective needs of a household, if given the option, households typically prefer larger housing units with more rooms and amenities in a better condition. A crucial aspect is the type of building and tenure (e.g., owner-occupied single-family home vs. rental apartment). The most significant part of income segregation can be attributed to a spatially unequal supply of housing of different types \citep{loufPatternsResidentialSegregation2016, oecdDividedCitiesUnderstanding2018, owensBuildingInequalityHousing2019}. I will subsume these properties of housing units under the label "housing quality." My first assumption, therefore, reads:

\begin{quotation}
\textbf{Assumption 1:} Households prefer a higher housing quality.
\end{quotation}

As already established, households primarily care about the housing quality provided by a unit if they consider to move there \citep{chauCriticalReviewLiterature2003, bayerEquilibriumModelSorting2004, mummoloWhyPartisansNot2017, delucaNotJustLateral2020}.


\subsubsection{Preferences for Neighbors and Neighborhoods}

One strand of the literature suggests that households prefer the homophily of a neighborhood. For example, households appear to prefer to live next to neighbors with similar education \citep{bayerUnifiedFrameworkMeasuring2007, bruchChoiceSetFormation2019}, race and ethnicity \citep{krysanDoesRaceMatter2009, bruchChoiceSetFormation2019}, as well as attitudes, beliefs and lifestyles \citep{mummoloWhyPartisansNot2017}. Especially when considering multiple characteristics simultaneously, homophily seems to be a well-fitting general principle in neighborhood formation \citep{loganBirdsFeatherSocial2016}. Such preferences might not just be about the relationships households might have with their neighbors, but living in a neighborhood with similar others results in businesses and organizations catering to the preferences and needs of these residents \citep{bruecknerWhyCentralParis1999}.

However, one could also assume that households prefer neighbor characteristics they associate with higher neighborhood quality. Households prefer rich neighbors, hoping this will help finance public goods in the neighborhood \citep{bayerUnifiedFrameworkMeasuring2007, bruchChoiceSetFormation2019}. This can be interpreted as a preference for neighborhood quality, e.g., low crime and high school quality \citep{bayerUnifiedFrameworkMeasuring2007, mummoloWhyPartisansNot2017}. The extensive literature on racial neighborhood preferences in the US also suggests that the observed preferences are not so much attributable to homophily but whites' outgroup avoidance and fears of declining property values and increased crime \citep{krysanDoesRaceMatter2009}. Whites also tend to prefer the presence of Asian Americans over Hispanics over Blacks, while blacks prefer integrated neighborhoods \citep{charlesDynamicsRacialResidential2003}. Racial neighborhood preferences reflect the status hierarchy of racial groups.

Although both assumptions are plausible, I decided to go for the latter assumption that households want to maximize the status of their neighbors as implemented by \citet{benardWealthStatusBasedModel2007}\footnote{During theory-building, I have experimented with both assumptions and was not able to implement a model assuming homophily that produced the reviewed stylized facts}. It is also likely that households prefer similarity in some characteristics while preferring higher or lower levels in others. I subsume the different aspects of social characteristics into a single measure of the household's social status, which can correlate with their income. This status variable indicates the desirability of the household as a neighbor. My second assumption states:

\begin{quotation}
\textbf{Assumption 2:} Households prefer a higher social status of their neighbors.
\end{quotation}

There is also some discussion about which scale of neighborhood is relevant to household decisions. \citet{loganBirdsFeatherSocial2016} report previous studies and present evidence suggesting households consider relatively small scales around their location most. 


\subsection{Supply: Uncertainty, Investments, and Inertia}

One of the peculiarities of the housing market is that housing supply is comparably \textit{inelastic}; supply can only adapt to demand slowly as building new housing takes time. Especially in large, densely populated cities, there is almost no possibility of creating additional housing units \citep{baum-snowConstraintsCityNeighborhood2023}. Additionally, relative to the existing number of housing units, newly added or demolished buildings only marginally change the total number of housing units in a city \citep[1569]{whiteheadUrbanHousingMarkets1999}. I will, therefore, model cities with a fixed number of housing units. The main activity of suppliers of housing (landlords) is \textit{maintenance} of \textit{existing} housing stock. They need to decide how much they invest into the quality of their property as it would otherwise decline \citep[119]{bruecknerLecturesUrbanEconomics2011}. Empirical analyses suggest that the values of housing units decrease by about 0.5 to 2.5 \% annually \citep[1086]{rosenthalChangePersistenceEconomic2015}. These results also show that even if landlords do not invest, the physical conditions of housing units change only slowly. The same holds true for upgrading a unit as well: planning and construction take time and the larger the investment, the longer construction will take. The two processes are asymetrical though. Physical decay will be slower than (re)development \citep{glaeserUrbanDeclineDurable2005}. Nonetheless, the housing quality at a particular location (independent from its neighborhood) can change over time through physical decay from disinvestment, investments into the existing housing stock or even demolition and new construction.

\begin{quotation}
\textbf{Assumption 3:} Housing quality is variable but path dependent and has a lot of inertia.
\end{quotation}

However, there is uncertainty about the neighborhood's future development, making investments in built infrastructure risky since they need quite a long time to pay off. The reason of the uncertainty of investment is that the landlord has not full control over the desirability of their unit as the neighborhood matters for households as well. The landlord is faced with fundamental uncertainty about an investment's returns and needs to rely on their predictions of the future of the area \citep{whiteheadUrbanHousingMarkets1999, galsterNatureNeighbourhood2001, glaeserExtrapolativeModelHouse2017}. Prices of other housing units are the most critical information investors consider when inferring demand in housing markets, although they are only available with a significant time lag \citep{glaeserExtrapolativeModelHouse2017}. I assume that landlords update their expectations about the neighborhood's future and the profitability of investments based on developments in their neighborhood \citep{galsterNatureNeighbourhood2001, bayerSpeculativeFeverInvestor2021}. 

\begin{quotation}
\textbf{Assumption 4:} Landlords invest in housing quality if the average rent in their neighborhood increases.
\end{quotation}

So, if a neighborhood becomes more attractive, either because of population changes in the neighborhood or because other landlords invested in their housing quality, competition for housing units in that neighborhood increases and rents will rise. As rents rise, landlords update their expectations about profitability and invest into their housing quality. If average rents rise only slightly, landlords will only invest marginally, but when rents rise fast, landlords make larger investments.

\subsection{Rent}

The price or rent the household pays the landlord to live at a given housing unit is a result of supply and demand in competitive markets. Specifically, rent is a mapping of desirability of the housing unit onto the income distribution. The housing units providing the highest residential satisfaction will have the fiercest competition between interested renters, and the renters with the highest disposable income can outprice the others. Therefore, the units with the highest residential satisfaction will be the most expensive, while the least expensive have the lowest level of residential satisfaction. To implement such a price-building mechanism, I decided to model rent as the 75th percentile of the incomes of all renters living in a housing unit with a lower residential satisfaction than their unit provides. Renters living in housing units with lower residential satisfaction want to move there because they can increase their residential satisfaction by moving to this housing unit. The renters in better housing units have no incentive to move to this unit. The 75th percentile is somewhat arbitrary but gives some leeway for movements to occur, whereas asking the maximum income of renters in a lower utility housing unit as rent makes it very unlikely that a household will move there. The lower rent than potentially achievable reduces the time the landlord needs to wait until a renter willing to pay the asking price moves in. %In general, renting markets are known to be comparably inefficient (CITE).